\documentclass[11pt,a4paper]{article}

\usepackage[utf8]{inputenc}
\usepackage[T1]{fontenc}
\usepackage{lmodern}
\usepackage[margin=25mm]{geometry}
\usepackage{booktabs}
\usepackage{tabularx}
\usepackage{array}
\usepackage{longtable}
\usepackage{listings}
\usepackage{enumitem}
\usepackage{amsmath}
\usepackage{microtype}
\usepackage{needspace}
\usepackage[hidelinks]{hyperref}

\setlength{\parindent}{0pt}
\setlength{\parskip}{0.5em}
\setlength{\emergencystretch}{3em}
\setlist[itemize]{leftmargin=1.6em}
\setlist[enumerate]{leftmargin=1.8em}
\renewcommand{\arraystretch}{1.18}
\newcolumntype{Y}{>{\raggedright\arraybackslash}X}

\lstset{
  basicstyle=\ttfamily\small,
  frame=single,
  breaklines=true,
  columns=fullflexible,
  keepspaces=true,
  showstringspaces=false,
  tabsize=2
}

\newcommand{\file}[1]{\texttt{\detokenize{#1}}}
\newcommand{\code}[1]{\texttt{\detokenize{#1}}}

\title{Flapping-Wing SBUS Controller\\
       \large Control, Telemetry, Wiring, and Safety Reference}
\author{Repository engineering documentation}
\date{27 July 2026}

\begin{document}

\maketitle

\begin{abstract}
This document defines the maintained wireless controller for a flapping-wing
robot. It begins with the system contract, follows the forward SBUS control
pipeline, then follows battery voltage through the complete return telemetry
pipeline. It also defines direct Python control through PySerial, failsafe and
recovery behavior, build instructions, and known limitations. The optional GUI
is a client of the same serial protocol; it is not required for automation. The
source code remains authoritative if implementation and documentation differ.
\end{abstract}

\tableofcontents
\clearpage

\section{System Summary}

The maintained system has four components:

\begin{enumerate}
  \item a Python program on the PC, or the optional supplied GUI;
  \item a USB-connected ESP32 transmitter;
  \item an ESP32-C3 SuperMini receiver on the robot; and
  \item a flight controller that accepts inverted SBUS.
\end{enumerate}

\begin{lstlisting}
Forward control:
Python / optional GUI -> USB serial -> ESP32 transmitter -> ESP-NOW
                      -> ESP32-C3 receiver -> GPIO4 SBUS -> flight controller

Reverse telemetry:
GPIO3 ADC -> ESP32-C3 receiver -> ESP-NOW -> ESP32 transmitter
          -> USB serial -> Python / optional GUI
\end{lstlisting}

The PC transmits six channel values in the exact order shown in
Table~\ref{tab:channels}. CH4 and CH7 are not part of the PC command.

\begin{table}[ht]
\centering
\caption{PC channel order and safe initial values.}
\label{tab:channels}
\begin{tabularx}{\textwidth}{@{}c c Y c Y@{}}
\toprule
Order & Channel & Function & Safe value & Meaning \\
\midrule
1 & CH1 & Yaw & 1500 & Centered \\
2 & CH2 & Pitch & 1500 & Centered \\
3 & CH3 & Throttle & 1000 & Minimum \\
4 & CH5 & Trim 1 / servo center & 1500 & Centered \\
5 & CH6 & Trim 2 / servo center & 1500 & Centered \\
6 & CH8 & Throttle lock / arm & 1000 & Locked; 2000 is armed \\
\bottomrule
\end{tabularx}
\end{table}

The complete host command is newline-terminated ASCII:

\begin{lstlisting}
<CH1,CH2,CH3,CH5,CH6,CH8>\n
\end{lstlisting}

For example, the safe locked command is:

\begin{lstlisting}
<1500,1500,1000,1500,1500,1000>\n
\end{lstlisting}

Host values are dimensionless control units from 1000 to 2000. They are not
PWM pulse widths in microseconds. The receiver maps the host range to raw SBUS
counts: 1000 to 172, 1500 to approximately 992, and 2000 to
1811~\cite{bolderflight-sbus}.

The normal physical connections are summarized below. The receiver USB port is
used for flashing and optional receiver diagnostics; normal PC control uses only
the transmitter USB port.

\begin{table}[ht]
\centering
\begin{tabularx}{\textwidth}{@{}lY@{}}
\toprule
Connection & Purpose \\
\midrule
PC to transmitter USB & Bidirectional PC serial transport and transmitter power \\
Receiver GPIO4 to flight-controller SBUS input & Inverted 3.3~V SBUS signal \\
Receiver GPIO3 to divider midpoint & Battery ADC measurement \\
Receiver GND to flight-controller GND and battery negative & Common reference \\
\bottomrule
\end{tabularx}
\end{table}

\section{SBUS Control Transmission Pipeline}

\subsection{Stage 1: Python to transmitter over USB serial}

The PC opens the transmitter COM port at 115200 baud, 8N1. In this notation,
\textbf{8} means eight data bits, \textbf{N} means no parity bit, and
\textbf{1} means one stop bit. An asynchronous UART frame also contains a start
bit. The 115200 baud rate specifies 115200 signaling intervals per second; for
this binary UART it is nominally 115200 bit/s. Each 8N1 character occupies 10
bits including the start and stop bits, so the theoretical maximum is about
11520 characters/s before USB and software overhead.

Commands are ASCII lines and every command must end with a newline byte
(\code{\n}, hexadecimal \code{0A}). The PC should first send
\code{IDENTIFY\n} and require this exact response before it sends channel data:

\begin{lstlisting}
DEVICE:FLAPPING_WING_TRANSMITTER;PROTOCOL:2;RADIO:1
\end{lstlisting}

This check reduces the risk of controlling the wrong COM-port device. A valid
channel line updates all six exposed controls at once. The transmitter does not
return a per-command acknowledgement containing the applied values. The PC
must continue sending a complete channel command at least every 100~ms; a gap
longer than 500~ms triggers the transmitter host failsafe.

\subsection{Stage 2: transmitter to receiver over ESP-NOW}

The transmitter maps the six PC values into their positions in a packed
\code{SbusPacket}. The structure contains sixteen unsigned 16-bit channel
slots and one 8-bit host-failsafe flag, for a total of 33 bytes. It is sent
immediately after a valid command; 50~ms wireless heartbeats retransmit the
latest state.

ESP-NOW uses 2.4~GHz Wi-Fi channel~1 and broadcast delivery. The current link is
not encrypted or authenticated. The binary structure is an internal firmware
interface, not the PC serial format. Its single authoritative definition is
\file{firmware/link/esp_now_link.h}; both MCUs must be reflashed after a
payload-format change.

\subsection{Stage 3: receiver qualification and SBUS output}

The ESP32-C3 receiver maps each host control value $x$ to raw SBUS counts:

\[
y = 172 + \operatorname{round}\left(\frac{(x-1000)(1811-172)}{1000}\right).
\]

After its startup and link-qualification sequence described in
Section~\ref{sec:failsafe}, the receiver enables UART1 and sends inverted SBUS
on GPIO4 approximately every 10~ms. The ESP32-C3 GPIO matrix permits peripheral
outputs such as UART TX to be routed to GPIO pins~\cite{espressif-c3-gpio}.

\begin{lstlisting}
ESP32-C3 GPIO4 ---------------- Flight-controller SBUS input
ESP32-C3 GND ------------------ Flight-controller GND
\end{lstlisting}

SBUS uses 100000 baud, 8E2: eight data bits, even parity, and two stop bits. The
electrical signal is inverted. These settings are intentionally different from
the PC USB-serial connection, which is 115200 baud, 8N1, and ordinary polarity.
The serial format, inversion, raw channel range, and typical update interval are
documented by the SBUS library used by this project~\cite{bolderflight-sbus}.

The embedded transmitter and receiver do not ramp or smooth channel updates.
Manual changes take effect immediately. Only the optional GUI's automated
benchmark ramps CH3, at 250 host units per second.

\subsection{Control timing summary}

\begin{table}[ht]
\centering
\begin{tabularx}{\textwidth}{@{}YYr@{}}
\toprule
Link & Representation & Normal interval \\
\midrule
Python to transmitter & Newline-ended ASCII & 100 ms recommended \\
Transmitter to receiver & 33-byte binary ESP-NOW packet & 50 ms heartbeat \\
Receiver to flight controller & Inverted SBUS, 100000 baud 8E2 & 10 ms \\
\bottomrule
\end{tabularx}
\end{table}

\section{Voltage Measurement and Return Telemetry Pipeline}

\subsection{Stage 1: battery divider to GPIO3}

The ESP32-C3 measures the divider midpoint on GPIO3, which is the
\code{ADC1_CH3} input~\cite{espressif-c3-datasheet}. The 300~k$\Omega$ upper
resistor and 100~k$\Omega$ lower resistor form a nominal 4:1 divider:

\[
V_{GPIO3}=V_{battery}\frac{100}{300+100}=\frac{V_{battery}}{4},
\qquad
V_{battery}=4V_{GPIO3}.
\]

\begin{lstlisting}
Battery + ---- 300 kOhm ----+-- ESP32-C3 GPIO3
                            |
                          100 kOhm
                            |
Battery - ------------------+-- ESP32-C3 GND
\end{lstlisting}

Battery negative, receiver ground, and flight-controller ground must share a
common reference. Never connect the battery directly to GPIO3. Power the board
only through a board-appropriate regulated input or USB; do not apply raw
battery voltage to an MCU supply pin unless the exact board and regulator are
rated for it.

The firmware configures 12-bit conversion and 11~dB attenuation. Espressif
documents an approximate 0--2500~mV measurable range for the ESP32-C3 at this
attenuation, so the nominal divider corresponds to approximately 0--10~V at the
battery~\cite{arduino-esp32-adc}. This is not a precision or absolute-maximum
guarantee. Preserve margin for resistor tolerance, ADC error, noise, and battery
transients, and compare the result with a multimeter over the intended range.

The divider has a 75~k$\Omega$ Thevenin source resistance. ADC sampling can be
sensitive to source impedance and noise. If readings are unstable or biased,
place a 100~nF bypass capacitor close to GPIO3 and ground, then calibrate the
complete divider rather than assuming the nominal ratio alone. Espressif
specifically recommends a 0.1~$\mu$F capacitor between an ESP32-C3 ADC pin and
ground and recommends use of ADC1~\cite{espressif-c3-hardware-guidelines}.

\subsection{Stage 2: receiver sampling}

Every telemetry update, the receiver averages 16 samples of both Arduino ADC
interfaces:

\begin{itemize}
  \item \code{battery_adc_raw}: averaged original 12-bit ADC count, nominally
        0--4095;
  \item \code{battery_pin_mv}: averaged calibrated GPIO3 voltage in
        millivolts~\cite{arduino-esp32-adc}.
\end{itemize}

The raw count is retained for diagnostics. The millivolt value is the preferred
input for converting the divider voltage into a displayed battery voltage.

\subsection{Stage 3: receiver to transmitter over ESP-NOW}

The receiver sends a packed \code{ReceiverStatusPacket} every 250~ms. Its exact
15-byte layout is shown below; multibyte fields use the ESP32's native
little-endian representation.

\begin{table}[ht]
\centering
\begin{tabularx}{\textwidth}{@{}lcrY@{}}
\toprule
Field & Type & Bytes & Meaning \\
\midrule
\code{magic} & uint32 & 4 & Constant \code{0x46575354} (FWST) \\
\code{packets_received} & uint32 & 4 & Accepted control-packet counter \\
\code{battery_adc_raw} & uint16 & 2 & Averaged raw ADC count \\
\code{battery_pin_mv} & uint16 & 2 & Averaged GPIO3 millivolts \\
\code{version} & uint8 & 1 & Receiver-status version, currently 2 \\
\code{link_active} & uint8 & 1 & 1 only while qualified SBUS output is active \\
\code{failsafe} & uint8 & 1 & 1 while the receiver reports a fault state \\
\bottomrule
\end{tabularx}
\end{table}

The receiver unicasts the status to the MAC address from which it most recently
accepted control. The transmitter checks the packet length, magic value, and
version before exposing it to the PC.

\subsection{Stage 4: transmitter ASCII line to Python}

The transmitter converts the binary status packet to a newline-ended ASCII line
on the same USB COM port used for commands:

\begin{lstlisting}
RECEIVER_STATUS mac=AA:BB:CC:DD:EE:FF packets=1234 link=1 failsafe=0 battery_raw=1800 battery_pin_mv=1100
\end{lstlisting}

This telemetry version reports receiver health and battery measurement but does
not return channel values or acknowledge individual PC commands.

The complete return pipeline is therefore:

\begin{lstlisting}
Battery voltage
  -> 300/100 kOhm divider
  -> GPIO3 ADC (16-sample average)
  -> binary ReceiverStatusPacket over ESP-NOW
  -> transmitter validates and formats ASCII
  -> USB serial receive buffer
  -> Python ser.readline()
  -> parse link, failsafe, and packet count
  -> parse battery_pin_mv
  -> multiply by 4 and divide by 1000
  -> battery volts
\end{lstlisting}

For example, \code{battery_pin_mv=1100} represents 1.100~V at GPIO3 and an
estimated battery voltage of $1.100\times4=4.40$~V. Battery percentage is not
estimated because a meaningful percentage requires battery chemistry, cell
count, load behavior, calibration, and chosen operating thresholds.

\section{Direct Python Control with PySerial}

\subsection{Install PySerial and locate the port}

The PyPI package is named \code{pyserial}, while the module imported by Python
is named \code{serial}:

\begin{lstlisting}
py -m pip install pyserial
py -m serial.tools.list_ports -v
\end{lstlisting}

Select the transmitter port by verified identity, not by assuming a fixed COM
number. Only one process can own a serial port at a time; close PlatformIO
Serial Monitor, Arduino Serial Monitor, and the GUI before opening the port in
another Python process.

\subsection{Opening, writing, and reading}

The explicit PySerial configuration for 115200 baud, 8N1, no flow control is:

\begin{lstlisting}[language=Python]
import serial

ser = serial.Serial(
    port="COM10",
    baudrate=115200,
    bytesize=serial.EIGHTBITS,
    parity=serial.PARITY_NONE,
    stopbits=serial.STOPBITS_ONE,
    timeout=0.05,
    write_timeout=0.25,
    xonxoff=False,
    rtscts=False,
    dsrdtr=False,
)
\end{lstlisting}

\code{ser.write(...)} requires bytes. Encode text commands as ASCII:

\begin{lstlisting}[language=Python]
values = [1500, 1500, 1000, 1500, 1500, 1000]
command = ("<" + ",".join(map(str, values)) + ">\n").encode("ascii")
ser.write(command)
\end{lstlisting}

Bytes received from the transmitter accumulate in the operating-system and
PySerial receive buffers. \code{ser.readline()} returns bytes through the next
newline, or returns the available bytes when the configured timeout expires.
Decode after reading:

\begin{lstlisting}[language=Python]
line = ser.readline().decode("ascii", errors="replace").strip()
\end{lstlisting}

For event-driven code, \code{ser.in_waiting} reports the bytes currently
available and \code{ser.read(n)} retrieves up to \code{n} bytes. Such code must
retain partial data between reads and split only at newline boundaries.

\subsection{Minimal safe bidirectional client}

The following program identifies the correct firmware, maintains a safe locked
control heartbeat, and parses telemetry from the same serial port. Opening the
port can reset an ESP32, so identification is retried for five seconds.

\begin{lstlisting}[language=Python]
import time
import serial

PORT = "COM10"
SAFE = [1500, 1500, 1000, 1500, 1500, 1000]
EXPECTED = "DEVICE:FLAPPING_WING_TRANSMITTER;PROTOCOL:2;RADIO:1"

def encode_channels(values):
    if len(values) != 6:
        raise ValueError("six channel values are required")
    return ("<" + ",".join(str(v) for v in values) + ">\n").encode("ascii")

with serial.Serial(
    port=PORT,
    baudrate=115200,
    bytesize=serial.EIGHTBITS,
    parity=serial.PARITY_NONE,
    stopbits=serial.STOPBITS_ONE,
    timeout=0.05,
    write_timeout=0.25,
    xonxoff=False,
    rtscts=False,
    dsrdtr=False,
) as ser:
    ser.reset_input_buffer()

    deadline = time.monotonic() + 5.0
    next_identify = 0.0
    while True:
        now = time.monotonic()
        if now >= deadline:
            raise RuntimeError("expected transmitter not found")
        if now >= next_identify:
            ser.write(b"IDENTIFY\n")
            next_identify = now + 0.5
        line = ser.readline().decode("ascii", errors="replace").strip()
        if line == EXPECTED:
            break

    safe_packet = encode_channels(SAFE)
    next_control = time.monotonic()
    try:
        while True:
            now = time.monotonic()
            if now >= next_control:
                ser.write(safe_packet)
                next_control = now + 0.100

            line = ser.readline().decode("ascii", errors="replace").strip()
            if line.startswith("RECEIVER_STATUS "):
                try:
                    fields = dict(part.split("=", 1)
                                  for part in line.split()[1:])
                    battery_v = int(fields["battery_pin_mv"]) * 4 / 1000
                except (ValueError, KeyError):
                    continue
                print(f"packets={fields['packets']} "
                      f"link={fields['link']} "
                      f"failsafe={fields['failsafe']} "
                      f"battery={battery_v:.2f} V")
    finally:
        for _ in range(3):
            try:
                ser.write(safe_packet)
            except serial.SerialException:
                break
            time.sleep(0.05)
\end{lstlisting}

A production client should validate six integers before encoding them, keep CH8
at 1000 until deliberate arming, tolerate ESP32 boot text and malformed lines,
track a single receiver MAC, detect stale telemetry, and explicitly send the
locked state before closing. The context manager closes the port even if an
exception occurs. Full message grammars and a more state-aware example are in
\file{docs/serial-protocol.md}.

\subsection{Live command and battery CSV recording}

Reusable parsing and recording functions live in \file{host/pc_tools/}.
The GUI calls the same \code{LiveCsvLogger} that is available to custom Python
programs:

\begin{lstlisting}[language=Python]
from host.pc_tools import LiveCsvLogger, parse_receiver_status

status = parse_receiver_status(line)
if status is not None:
    logger.log_receiver_status(status, commanded_channels)
\end{lstlisting}

Run the GUI and use its \textbf{START RECORDING} button:

\begin{lstlisting}
py -m host.gui
\end{lstlisting}

For safe locked-state recording without the GUI:

\begin{lstlisting}
py -m host.log_telemetry --list-ports
py -m host.log_telemetry --port COM10 --duration 60
\end{lstlisting}

Each CSV status row contains PC time, commanded channels, link/failsafe state,
packet count, raw ADC, GPIO3 millivolts, and reconstructed battery volts.
Commanded channels are PC-side records, not receiver acknowledgements. Rows are
flushed immediately. The GUI and standalone recorder cannot own the same COM
port simultaneously.

\subsection{PC serial message reference}

\begin{table}[ht]
\centering
\begin{tabularx}{\textwidth}{@{}lYY@{}}
\toprule
PC sends & Response type & Purpose \\
\midrule
\code{IDENTIFY\n} & \code{DEVICE:} identity line & Verify firmware, protocol 2, and radio initialization \\
\code{PING\n} & \code{PONG;} health line & Query failsafe, radio, and uptime state \\
Six-channel line & No command-specific response & Update all exposed controls \\
\bottomrule
\end{tabularx}
\end{table}

\code{RECEIVER_STATUS} lines are asynchronous: they arrive when telemetry is
received, not as replies to a particular PC command. They do not acknowledge
the receiver's applied channel state.

\section{Failsafe and Recovery}
\label{sec:failsafe}

\subsection{Receiver startup qualification}

\begin{enumerate}
  \item For five seconds after receiver boot, UART1 is disabled, GPIO4 is
        configured as an input (high-impedance), and no SBUS frames are sent.
  \item After the delay, the receiver requires five consecutive healthy packets
        from one transmitter with the host failsafe clear and CH8 locked at
        1000.
  \item It then enables SBUS and locks to that transmitter MAC until receiver
        reboot. A different transmitter cannot take over during that boot.
\end{enumerate}

\Needspace{0.42\textheight}
\subsection{Fault response}

\begin{table}[ht]
\centering
\begin{tabularx}{\textwidth}{@{}YY@{}}
\toprule
Condition & Implemented response \\
\midrule
No valid PC channel command for more than 500~ms & Transmitter preserves the latest CH3, forces CH8 to 1000, and sets its ESP-NOW host-failsafe flag \\
Receiver sees the host-failsafe flag & Receiver stops UART1 and makes GPIO4 high-impedance \\
Receiver gets no control packet for more than 500~ms & Receiver stops UART1 and makes GPIO4 high-impedance \\
\bottomrule
\end{tabularx}
\end{table}

The receiver does not send an SBUS frame with the SBUS protocol failsafe bit
set. It represents link loss by stopping SBUS output because the connected
flight controller's missing-SBUS behavior is the verified response.

\subsection{Recovery and safe operation}

Recovery cannot automatically arm the robot. The PC must first restore valid
commands with CH8=1000. The receiver requires five healthy locked packets before
it restores SBUS; only then may the PC deliberately change CH8.

Before closing the transmitter COM port, send the safe locked command several
times. If the process exits unexpectedly, the 500~ms host timeout still causes
the transmitter and receiver failsafe sequence.

Before first powered use after wiring or firmware changes:

\begin{enumerate}
  \item verify GPIO4 connects only to the flight-controller SBUS input;
  \item verify the common ground and the 300/100~k$\Omega$ divider orientation;
  \item confirm divided voltage with a multimeter before connecting GPIO3;
  \item restrain the mechanism and keep CH8 locked;
  \item verify identity, one receiver MAC, \code{link=1},
        \code{failsafe=0}, and plausible battery voltage; and
  \item test USB loss and radio loss before unrestricted operation.
\end{enumerate}

\section{Build, Flash, and Run}

The transmitter and receiver are independent PlatformIO projects. Replace the
example COM numbers with the ports reported on the current PC. Protocol~3
changes both packed ESP-NOW structures; flash both nodes before running the
updated GUI or logger. Mixed protocol firmware is rejected.

\subsection{Transmitter}

\begin{lstlisting}
cd firmware\transmitter
& "$env:USERPROFILE\.platformio\penv\Scripts\platformio.exe" run --target upload --upload-port COM10
cd ..\..
\end{lstlisting}

PlatformIO environment: \code{esp32doit-devkit-v1}, for an ESP32 Dev Module or
compatible DevKit board.

\subsection{Receiver}

\begin{lstlisting}
cd firmware\receiver
& "$env:USERPROFILE\.platformio\penv\Scripts\platformio.exe" run --target upload --upload-port COM11
cd ..\..
\end{lstlisting}

PlatformIO environment: \code{esp32-c3-devkitm-1}, used for the ESP32-C3
SuperMini. In Arduino IDE, select \textbf{ESP32C3 Dev Module}. Enable
\textbf{USB CDC On Boot} only when native-USB receiver debug output is required;
it does not create the SBUS link or the operational PC control connection.

\subsection{One-board wired ESP32-C3 diagnostic}

The separate \file{diagnostics/wired_c3/} project bypasses ESP-NOW so GPIO4
SBUS wiring and the GPIO3 battery-divider input can be tested with one C3. Power
the C3 from PC USB, connect GPIO4 to the flight-controller SBUS input, and join
the C3 ground, flight-controller ground, and battery negative. Leave the
flight-controller 5~V wire disconnected while USB powers the board. GPIO3 may
remain connected to the 300/100~k$\Omega$ divider midpoint.

\begin{lstlisting}
cd diagnostics\wired_c3
& "$env:USERPROFILE\.platformio\penv\Scripts\platformio.exe" run --target upload --upload-port COM11
cd ..\..
py -m diagnostics.wired_c3.gui
\end{lstlisting}

The diagnostic preserves the five-second silent startup, requires repeated
locked commands before enabling SBUS, stops GPIO4 output after a 500~ms host
timeout, and reports GPIO3 raw ADC and voltage to its dedicated GUI. Flash
\file{firmware/receiver/} again to restore wireless operation after the test.

List available ports with:

\begin{lstlisting}
& "$env:USERPROFILE\.platformio\penv\Scripts\platformio.exe" device list
\end{lstlisting}

PlatformIO may auto-detect a port when only one compatible board is connected,
but explicit \code{--upload-port} avoids flashing the wrong MCU. Stop any serial
monitor or Python client before upload because the COM port cannot normally be
opened by two processes simultaneously.

\section{Optional GUI}

Install its dependencies and run it from the repository root:

\begin{lstlisting}
py -m pip install -r requirements.txt
py -m host.gui
\end{lstlisting}

The GUI identifies the transmitter before enabling control, monitors radio and
receiver telemetry, displays reconstructed battery voltage, locks CH8 on
detected faults, and requires confirmation for arming and benchmark start.
Enter a CSV filename and select
\textbf{START RECORDING}; the file is created under \file{logs/}, the
\file{.csv} extension is added when omitted, and existing files are not
overwritten. It implements the same serial protocol described in Section~4.

Manual slider changes are transmitted immediately when the GUI applies the
change. Only the automated benchmark ramps CH3, at 250 host units per second.
The flashed \file{.ino} firmware performs no ramping.

\section{Repository, Development Record, and Future Work}

\subsection{Repository structure}

\begin{lstlisting}
diagnostics/wired_c3/          One-board USB-to-SBUS and GPIO3 test
firmware/link/                 ESP-NOW link packet definitions
firmware/transmitter/          ESP32 transmitter firmware/PlatformIO project
firmware/receiver/             ESP32-C3 receiver firmware/PlatformIO project
host/gui.py                    GUI, benchmark, and CSV recording entry point
host/log_telemetry.py          Standalone safe-state CSV recorder
host/pc_tools/                 Shared serial protocol and CSV logger
tests/test_pc_tools.py          PC protocol/logger unit tests
logs/                           Local CSV output; ignored by Git
docs/serial-protocol.md        Direct Python automation protocol
docs/flapping-wing-sbus-controller.tex
                               This engineering document
requirements.txt               Python runtime dependencies
pyproject.toml                 Python project metadata and entry points
output/pdf/                    Generated documentation PDF; ignored by Git
\end{lstlisting}

\code{RECEIVER_VALUE_TEST} in \file{firmware/receiver/receiver.ino} is disabled
by default. Enable it only for receiver-side serial diagnostics.

\subsection{Development record}

The original wired USB-to-SBUS prototype was split into two ESP32 nodes joined
by ESP-NOW. The robot-side node was moved to an ESP32-C3 SuperMini and its SBUS
output assigned to GPIO4. Subsequent work added host and radio failsafes,
locked-link startup qualification, receiver-to-transmitter telemetry, GPIO3
battery measurement through the 300/100~k$\Omega$ divider, direct Python
automation, and optional GUI safety handling.

\subsection{Known limitations and future directions}

Current limitations are:

\begin{itemize}
  \item ESP-NOW broadcast control is unauthenticated and unencrypted;
  \item both nodes use a fixed Wi-Fi channel;
  \item the forward packet has no explicit packet-version field;
  \item 4~Hz telemetry reports receiver health and battery data but does not
        acknowledge applied channels;
  \item battery measurement has no stored per-unit calibration or
        chemistry-aware percentage; and
  \item automated tests do not yet exercise the full serial, wireless, and
        failsafe state machines on hardware.
\end{itemize}

Prioritized improvements are explicit forward-packet versioning,
authenticated peers, configurable channel/peer provisioning,
calibrated voltage thresholds, telemetry freshness/quality metrics, and
repeatable hardware-in-the-loop disconnect and recovery tests.

\section{External References}

\begin{thebibliography}{9}

\bibitem{espressif-c3-datasheet}
Espressif Systems,
\href{https://documentation.espressif.com/esp32-c3_datasheet_en.pdf}
     {\emph{ESP32-C3 Series Datasheet}, version 2.4}.

\bibitem{espressif-c3-gpio}
Espressif Systems,
\href{https://docs.espressif.com/projects/esp-idf/en/v5.3.2/esp32c3/api-reference/peripherals/gpio.html}
     {\emph{GPIO \& RTC GPIO --- ESP32-C3}},
ESP-IDF Programming Guide, version 5.3.2.

\bibitem{espressif-c3-hardware-guidelines}
Espressif Systems,
\href{https://docs.espressif.com/projects/esp-hardware-design-guidelines/en/latest/esp32c3/schematic-checklist.html\#adc}
     {\emph{ESP32-C3 Hardware Design Guidelines: Schematic Checklist --- ADC}}.

\bibitem{arduino-esp32-adc}
Espressif Systems,
\href{https://docs.espressif.com/projects/arduino-esp32/en/latest/api/adc.html}
     {\emph{Arduino-ESP32 ADC API}}.

\bibitem{bolderflight-sbus}
Bolder Flight Systems,
\href{https://github.com/bolderflight/sbus\#description}
     {\emph{SBUS Library: Protocol Description and Inverted Serial Support}}.

\bibitem{pyserial-api}
PySerial contributors,
\href{https://pyserial.readthedocs.io/en/stable/pyserial.html}
     {\emph{PySerial API Documentation}}.

\bibitem{pyserial-tools}
PySerial contributors,
\href{https://pyserial.readthedocs.io/en/latest/tools.html}
     {\emph{PySerial Tools Documentation}}.

\bibitem{espressif-c3-adc-calibration}
Espressif Systems,
\href{https://docs.espressif.com/projects/esp-idf/en/latest/esp32c3/api-reference/peripherals/adc/adc_calibration.html}
     {\emph{ESP32-C3 ADC Calibration Driver}}.

\end{thebibliography}

\clearpage
\appendix
\section{Protocol Quick Reference}

\begin{longtable}{@{}p{0.29\textwidth}p{0.65\textwidth}@{}}
\toprule
Item & Value \\
\midrule
\endhead
PC serial port & Transmitter USB COM port, 115200 baud, 8N1 \\
Identity request & \code{IDENTIFY\n} \\
Expected identity & \code{DEVICE:FLAPPING_WING_TRANSMITTER;PROTOCOL:2;RADIO:1} \\
Control line & \code{<CH1,CH2,CH3,CH5,CH6,CH8>\n} \\
Safe line & \code{<1500,1500,1000,1500,1500,1000>\n} \\
PC heartbeat & At least one valid control line per 500~ms; 100~ms recommended \\
ESP-NOW control & Channel 1, 33-byte packed packet, 50~ms heartbeat \\
SBUS output & GPIO4, inverted, 100000 baud 8E2, approximately 10~ms \\
Battery input & GPIO3 through 300/100~k$\Omega$ divider \\
Telemetry & 15-byte battery/link status every 250~ms \\
Battery conversion & \code{battery_pin_mv / 1000 * 4} volts \\
Receiver fault action & Stop UART1 and make GPIO4 high-impedance \\
\bottomrule
\end{longtable}

\end{document}
